\clearpage
\VerdierPhysicalPageAnchor{0147}%
\begingroup
\small
\setlength{\parskip}{0pt}%
\setlength{\emergencystretch}{1.5em}%
\setlength{\abovedisplayskip}{3pt plus 1pt minus 1pt}%
\setlength{\belowdisplayskip}{3pt plus 1pt minus 1pt}%

\VerdierSemanticAnchor{ch2-2-3-4-d-complete}
\noindent 这证明了 d)；现在证明 f)。
\VerdierSemanticAnchor{ch2-2-3-4-f}
$\mathcal D$ 的对象 $Z$ 属于 $Q^{-1}(\operatorname{Lib}_R(Q))$，当且仅当
存在 $Q$-右自由对象 $Y$ 与 $\mathcal D/\mathcal B$ 中的同构
$Q(Z)\xrightarrow[\sim]{m}Q(Y)$。
\VerdierSemanticAnchor{ch2-2-3-4-f-fraction}
这种同构形如 $m=Q(s)^{-1}Q(t)$，其中
$s\in S_{\mathcal D}(\mathcal B)$ 且
$t\in\widehat{S_{\mathcal D}(\mathcal B)}$（(2.2.2) 与 (2.2.4)）。
因此在 $\mathcal D$ 中有图式：
\VerdierSemanticAnchor{ch2-2-3-4-f-diagram}
\begin{center}
  \begin{tikzcd}[column sep=4.1em,row sep=2.5em]
    Z \arrow[r,"t"] & X & Y \arrow[l,"s"']
  \end{tikzcd}\quad .
\end{center}
\VerdierSemanticAnchor{ch2-2-3-4-f-retraction}
由 a), iv)，态射 $s$ 有收缩 $s'\colon X\longrightarrow Y$，且该收缩
属于 $S_{\mathcal D}(\mathcal B)$ (1.2.6)。因此存在态射
$t'\colon Z\longrightarrow Y$，$t'\in\widehat{S_{\mathcal D}(\mathcal B)}$。
\VerdierSemanticAnchor{ch2-2-3-4-f-converse}
反之，若存在这样的态射，则 $Q(t')\colon Q(Z)\longrightarrow Q(Y)$ 是同构。
\VerdierSemanticAnchor{ch2-2-3-4-g}
还需证明 g)。由 a), v)，有双函子同构：
\VerdierSemanticAnchor{ch2-2-3-4-g-q1-adjunction}
\[
 \operatorname{Hom}_{Q^{-1}(\operatorname{Lib}_R(Q))}(Z,\Psi(W))
 \xrightarrow{\sim}
 \operatorname{Hom}_{\operatorname{Lib}_R(Q)}(Q_1(Z),W),
\]
这表明 $\Psi$ 是 $Q_1$ 的右伴随。
\VerdierSemanticAnchor{ch2-2-3-4-g-inclusion-adjunction}
因此，若以 $i\colon L_R(Q)\lhook\joinrel\longrightarrow
Q^{-1}(\operatorname{Lib}_R(Q))$ 表示包含函子，则有双函子同构：
\[
 \operatorname{Hom}_{Q^{-1}(\operatorname{Lib}_R(Q))}(Z,i(Y))
 \xrightarrow{\sim}
 \operatorname{Hom}_{L_R(Q)}(\Psi Q_1(Z),Y)\, .
\]
\VerdierSemanticAnchor{ch2-2-3-4-g-proof-complete}
这表明 $\mathcal R=\Psi Q_1$ 是包含函子的左伴随，并完成命题 (2.3.3)
的证明。

\begingroup\itshape
\VerdierSemanticAnchor{ch2-2-3-5}
\medskip\noindent\textbf{命题 2.3.5.}\quad
\VerdierSemanticAnchor{ch2-2-3-5-setup}
设 $\mathcal D$ 为三角范畴，$\mathcal A$ 与 $\mathcal B$ 为两个满三角子范畴
（$\mathcal B$ 严格满），并设 $S_{\mathcal D}(\mathcal B)$ 与
$S_{\mathcal A}(\mathcal A\cap\mathcal B)$ 分别为 $\mathcal D$ 与
$\mathcal A$ 中与 $\mathcal B$ 和 $\mathcal A\cap\mathcal B$ 关联的乘法系。
\VerdierSemanticAnchor{ch2-2-3-5-a}
\par\noindent a) 下列条件等价：
\VerdierSemanticAnchor{ch2-2-3-5-a-i}
\par\noindent\hspace{1.2em}i) 每个态射 $X\xrightarrow{s}X'$，其中
$X\in\operatorname{Ob}(\mathcal A)$ 且 $s\in S_{\mathcal D}(\mathcal B)$，
都可嵌入交换图：
\VerdierSemanticAnchor{ch2-2-3-5-a-i-diagram}
\begin{center}
 \begin{tikzcd}[column sep=4.1em,row sep=2.6em]
 X \arrow[r,"s"] \arrow[dr,"{s'}"'] & X' \arrow[d] &
 {s'\in S_{\mathcal A}(\mathcal A\cap\mathcal B)} \\
 & X'' &
 \end{tikzcd}\quad .
\end{center}
\endgroup
\endgroup
